\subsection{最终结构的例子}

I. 结构的直接像。当 I 是由单个元素组成的集合时，关于 \((\mathbf{A}, \mathcal{S}, f)\) 的最终结构称为结构 G 在 f 下的直接像（当它存在时）。

\begin{enumerate}
\def\labelenumi{\Roman{enumi}.}
\setcounter{enumi}{1}
\tightlist
\item
  商结构。设 A 是赋予物种 \(\Sigma\) 的结构 S 的集合，设 R 是 A 上的等价关系，并设 \(\varphi\) 设φ为A到商集E = A/R上的典范映射（第二章，§6，第2节）。结构S在映射φ下的直接像（当它存在时）称为结构S关于关系R的商结构。
\end{enumerate}

* 一般而言，序结构或代数结构并不对任意等价关系都容许商结构（参见第三章第1节习题2）。另一方面，拓扑总是对任意等价关系容许商结构，但对于豪斯多夫拓扑则未必如此。 *

设 A、B 为两个集合，分别赋予结构 \(\mathcal{S}\) , \(\mathcal{S}'\) 的物种 \(\Sigma\) ，并且让 \(f\) 是A到B的一个态射。设R为等价关系 \(f(x) = f(y)\) ，让 \(\varphi\) 令 为 A 到 A/R 的典范映射，并令 \(j\) 为 \(f(A)\) 到B中。假设 \(\mathcal{S}\) 承认一个商结构 \(\mathcal{S}_0\) 关于R，并且 \(\mathcal{S}'\) 诱导一个结构 \(\mathcal{S}_0'\) 在 \(f(A)\) 上。那么，在典范分解 \(f = j \circ g \circ \varphi\) 的 \(f\) （第二章，§6，第5小节）中，由 \(g\) 给出的、把A/R映到 \(f(A)\) 上且与 \(f\) 相关联的双射，当A/R被赋予 \(\mathcal{S}_0\) 并且 \(f(A)\) 被赋予 \(\mathcal{S}_0'\) 时，是一个态射（但不一定是同构）。因为 \(j \circ g\) 根据商结构的定义是A/R到B的态射，因此 \(g\) 根据诱导结构的定义是A/R到 \(f(A)\)

的态射。 \(\mathcal{S}\) , \(\mathcal{S}'\) 的物种 \(\Sigma\) CST20. 设A，A'为两个带有结构 \(\mathcal{S}_0\) （分别） \(\mathcal{S}_0'\) )的 \(\mathcal{S}\) 的集合，并设R（相应地R'）为A（相应地A'）上的一个等价关系。假设存在一个由R（相应地由R'）得到的商结构 \(\mathcal{S}'\) 。如果 \(f\) 是A到A'的一个态射，且与等价关系R和R'相容，并且 \(g\) 是由 \(f\) 通过取商得到的映射，那么 \(g\) 是A/R到A'/R'的一个态射。

让 \(\varphi\) （分别） \(\varphi'\) ）是A到A/R（相应地A'到A'/R'）的典范映射；于是我们有 \(g \circ \varphi = \varphi' \circ f\) 。由于 \(\varphi'\) 并且 \(f\) 是态射，所以 \(\varphi' \circ f\) 由（MO \(_{II}\) ）。但接着， \(g \circ \varphi\) 作为态射， \(g\) 根据商结构的定义也是态射。

¶ 传递性判据CST19特别地给出以下判据：

CST21. 设A为带有结构 \(\mathcal{S}\) 的物种 \(\Sigma\) 的集合，并设R为A上的一个等价关系，使得在A/R上存在由R得到的商结构 \(\mathcal{S}'\) 的 \(\mathcal{S}\) 。设S为A上的一个比R更粗的等价关系，并设S/R表示A/R上由S得到的商等价关系（第二章，§6，第7小节）。那么存在(A/R)/(S/R)上由S得到的商结构 \(\mathcal{S}''\) 的 \(\mathcal{S}'\) 当且仅当在 A/S 上存在一个商结构 \(\mathcal{S}_0\) 的 \(\mathcal{S}\) ，并且A/S（带有 \(\mathcal{S}_0\) ）到(A/R)/(S/R)（带有 \(\mathcal{S}''\) ）的典范映射是一个同构。

让 \(\varphi\) 令 为 A 到 A/R 的典范映射，并令 \(\psi\) 为A/R到(A/R)/(S/R)的映射。根据CST19， \(\mathcal{G}''\) 是……的商 \(\mathcal{G}'\) 通过S/R当且仅当 \(\mathcal{G}''\) 是关于 (A, \(\mathcal{G}\) , \(\psi \circ \varphi\) ) 的最终结构。该准则随后由以下事实得出：关系 \(\psi(\varphi(x)) = \psi(\varphi(y))\) 等价于 \(S\{x,y\}\).

注记。设 A 为赋予了一个结构的集合 \(\mathcal{G}\) 的物种 \(\Sigma\) ，并设 R 是 A 上的一个等价关系，使得在 E = A/R 上存在一个商结构 \(\mathcal{G}'\) 的 \(\mathcal{G}\) 由 R 给出。设 \(\varphi\) 为 A 到 E 的典范映射。一般而言，不存在截面 \(s\) 的 \(\varphi\) （第二章，§3，第 8 号）是 E 到 A 的态射。假设这样的截面 \(s\) 存在，并且进一步假设存在一个结构 \(\mathcal{G}''\) ，由 \(\mathcal{G}\) 在 \(s(E)\) 上诱导。那么，若 \(j\) 表示 \(s(E)\) 到 A 的典范单射，且若 \(s = j \circ f\) ，双射 \(f\) 是 E 到 \(s(E)\) 的一个同构。因为 \(f\) 由诱导结构的定义是一个态射，而 \(g = \varphi \circ j\) 是……的一个态射 \(s(E)\) 到 E 的映射由于 \((MO_{II})\) 。由于 \(g \circ f\) 并且 \(f \circ g\) 分别是 E 和 \(s(E)\) 的恒等映射，该断言是 CST8 的推论。

